\section{Самосопряженные преобразования}
Пусть $\mathcal{E}$ -- евклидово пространство, причем $1\leqslant\dim \mathcal{E}=n<\infty$.

\begin{definition}
    Линейное преобразование $\phi$ называется самосопряженным, если 
    \[
    \forall x,y\in \mathcal{E}\hookrightarrow (\phi(x), y)=(x,\phi(y))
    \]
\end{definition}
Из определения и следствия теоремы 12.1 получаем, что преобразование $\phi$ самосопряженное тогда и только тогда, когда в ОНБ матрица $A_{\phi}$ симметричная (т.е. $A_{\phi}=A_{\phi}^T$).

\begin{theorem}
    Все корни характеристического многочлена самосопряженного преобразования вещественны.
\end{theorem}
\begin{proof}
    Докажем от противного. Пусть существует корень $\lambda\in \mathbb{C}$ характеристического многочлена $P(\lambda)$ преобразования $\phi$. Тогда по теореме 6.5 паре корней $\lambda, \overline{\lambda}$ будет соответствовать инвариантное двумерное подпространство $\mathcal{E}_1$, не содержащее собственных векторов.

    Рассмотрим ограничение преобразования $\tilde{\phi}$ на этом подпространстве. Оно самосопряжено. Выберем в $\mathcal{E}_1$ ОНБ. Тогда в нем матрица $A_{\tilde{\phi}}$ имеет симметричный вид
    \[
    A_{\tilde{\phi}}=\begin{pmatrix}
        \alpha & \beta \\
        \beta & \gamma
    \end{pmatrix}
    \]
    Тогда характеристический многочлен $\tilde{P}(\lambda)=(\alpha-\lambda)(\gamma-\lambda) - \beta^2$ имеет вещественные корни, так как $\tilde{P}(\alpha)=-\beta^2 \leqslant0$. Значит в $\mathcal{E}_1$ есть собственный вектор. Противоречие. Значит все корни $P(\lambda)$ вещественные.
\end{proof}

\textit{Следствие}: у любой симметричной матрицы все корни характеристического многочлена вещественные, так как в ОНБ можно рассмотреть самосопряженное преобразование с такой матрицей.

\begin{theorem}
    Собственные подпространства самосопряженного преобразования попарно ортогональны.
\end{theorem}
\begin{theorem}[\textbf{Эквивалентная формулировка теоремы 13.2}]
    Собственные векторы самосопряженного преобразования, отвечающие различным собственным значениям, ортогональны.
\end{theorem}
\begin{proof}
    Рассмотрим два неравных собственных значения $\lambda_1, \lambda_2$ и соответствующие им собственные векторы $h_1, h_2$. Из самосопряженности преобразования $\phi$
    \[
    (\phi(h_1), h_2)=(\lambda_1h_1,h_2), \quad (h_1, \phi(h_2))=(h_1,\lambda_2h_2) \Rightarrow (\lambda_1h_1,h_2)=(h_1,\lambda_2h_2)
    \]
    Значит $\lambda_1(h_1,h_2)=\lambda_2(h_1,h_2)$, откуда $(\lambda_1-\lambda_2)(h_1,h_2)=0$. Но $\lambda_1\neq\lambda_2$, следовательно $(h_1,h_2)=0$. Получили, что $h_1 \bot h_2$.
\end{proof}
\begin{theorem}
    Пусть $\mathcal{E}_1$ -- инвариантное подпространство самосопряженного преобразования $\phi$. Тогда ортогональное дополнение $\mathcal{E}_1^{\bot}$ тоже инвариантно относительно $\phi$.
\end{theorem}
\begin{proof}
    По определению $\forall x\in \mathcal{E}_1\hookrightarrow \phi(x)\in \mathcal{E}_1$. Значит $\forall y\in \mathcal{E}_1^{\bot} \hookrightarrow (\phi(x), y)=0$. Но из самосопряженности $\phi$ следует, что $(\phi(x), y)=(x,\phi(y))=0$, откуда получаем $\phi(y)\in \mathcal{E}_1^{\bot}$. Следовательно, $\mathcal{E}_1^{\bot}$ инвариантно относительно $\phi$.
\end{proof}
Теперь соберем из предыдущих трех теорем главную теорему.
\begin{theorem}
    Преобразование $\phi$ евклидова пространства $\mathcal{E}$ является самосопряженным тогда и только тогда, когда в $\mathcal{E}$ существует ОНБ из собственных векторов $\phi$.
\end{theorem}
\begin{proof}
    Если в $\mathcal{E}$ существует ОНБ из собственных векторов $\phi$, то в нем матрица преобразования имеет диагональный вид, а значит она симметричная, следовательно $\phi$ самосопряженное.

    Если $\phi$ -- самосопряженное, тогда все корни характеристического многочлена вещественные, а значит существует хотя бы одно собственное подпространство. Пусть $L$ -- прямая сумма всех собственных подпространств. Если $x\in L$, тогда $\phi(x)\in L$, так как $x$ и $\phi(x)$ -- линейные комбинации собственных векторов.

    Покажем, что $L\equiv \mathcal{E}$: предположим обратное.

    По теореме 13.4 из инвариантности $L$ следует инвариантность $L^{\bot}$, при этом по построению в $L^{\bot}$ нет собственных векторов. Рассмотрим ограничение ${\tilde\phi}$ на $L^{\bot}$. Так как $\tilde{\phi}$ -- самосопряженное, то по теореме 13.1 $\tilde{P}(\lambda)$ имеет вещественные корни, а значит в $L^{\bot}$ есть собственные векторы, противоречие.

    Значит $L\equiv \mathcal{E}$. Тогда выберем ОНБ в каждом слагаемом прямой суммы $L$. Их объединение будет базисом в $\mathcal{E}$, а по теореме 13.3 (с учетом единичной длины по построению) ортонормированным базисом.
\end{proof}
\textit{Следствие}: пусть $A$ -- симметричная матрица, тогда существует ортогональная матрица $S$ така, что $S^{-1}AS$ -- диагональная (рассматриваем ОНБ в котором $A$ симметричная и ОНБ из собственных векторов преобразования $A$; матрица перехода $S$ между этими базисами -- искомая).

\section{Ортогональные преобразования}
\begin{definition}
    Линейное преобразование $\phi$ евклидова пространства называется ортогональным, если
    \[
    \forall x,y\in \mathcal{E} \hookrightarrow (\phi(x),\phi(y))=(x,y)
    \]
\end{definition}

Из определения понятно, что $\phi$ сохраняет расстояния и углы. А значит $\phi$ биективно (из того, что $\ker \phi-\{0\}$, так как если $x\in\ker \phi$, то $(\phi(x),\phi(x))=(x,x)=0$ и $x=0$).

\begin{theorem}
    $\phi$ является ортогональным тогда и только тогда, когда $\phi^*=\phi^{-1}$.
\end{theorem}
\begin{proof}
    Если $\phi$ -- ортогональное, то $(x,y)=(\phi(x),\phi(y)=(x,\phi^*(\phi(y)))$. Значит $0=(x, \phi^*(\phi(y))-y)$, откуда по утверждению 10.1 $\phi^*(\phi(y))-y=0$, следовательно $\phi^*=\phi^{-1}$. Проделав выкладки в обратном порядке получаем доказательство в другую сторону.
\end{proof}
\begin{theorem}
    $\phi$ является ортогональным тогда и только тогда, когда его матрица $A$ ортогональна в любом ОНБ.
\end{theorem}
\begin{proof}
    По следствию из теоремы 12.1 в ОНБ $A^*=A^T$. Если $\phi$ ортогональное, то по теореме 14.1 $A^*=A^{-1}$, а значит $(A^*)^{-1}=A$ и $A^*(A^*)^{-1}=E=A^TA$. Если $A$ ортогональная, то $AA^T=E$, а значит $A^*=A^T=A^{-1}$ и по теореме 14.1 $\phi$ -- ортогональное.
\end{proof}

Если $e,f$ -- два ортонормированных базиса, тогда преобразование $\phi$: $\phi(e_i)=f_i, \forall i$ ортогонально по утверждению 11.2.
\begin{proposition}
    Если $\lambda \in \R$ -- собственное значение ортогонального преобразования $\phi$, то $|\lambda|=1$.
\end{proposition}
\begin{proof}
    Пусть $h$ -- соответствующий собственный вектор. Тогда 
    \[
    (h,h)=(\phi(h), \phi(h))=(\lambda h,\lambda h)=\lambda^2(h,h)\Rightarrow (\lambda^2-1)(h,h)=0
    \]
    Но так как по определению $h\neq 0$, то $(h,h)\neq0$ и $\lambda^2=1$.
\end{proof}
\textit{Замечание}: это утверждение верно и для $\lambda\in\mathbb{C}$.